% =========================================================
% Proof: Preservation of Diameter (Isometries)
% Mode: standard (stub)
% =========================================================

\subsection*{Preservation of Diameter}
\label{prf:preservation-of-diameter}

\begin{proof}
\Claim Let $(X,d_X)$ and $(Y,d_Y)$ be metric spaces and let
$f:X\to Y$ be an isometry. Then for every $A\subseteq X$,
\[
\operatorname{diam}(f(A))=\operatorname{diam}(A).
\]

\Given Metric spaces $(X,d_X)$, $(Y,d_Y)$, an isometry $f:X\to Y$, and a subset $A\subseteq X$.

\Goal To show $\operatorname{diam}(f(A))=\operatorname{diam}(A)$.

\Strategy Unfold the definition of diameter as a supremum of pairwise distances.
Use the defining property of an isometry to relate distances in $A$ to distances in $f(A)$.
Establish the needed set-inclusion(s) between the corresponding distance-sets and conclude equality of suprema.


Assume $f$ is an isometry, so
\[
d_Y\bigl(f(x),f(y)\bigr)=d_X(x,y)
\quad
\text{for all } x,y\in X.
\]

Let $x,y\in A$ be arbitrary.

Then the distance between their images equals the original distance:
\[
d_Y\bigl(f(x),f(y)\bigr)=d_X(x,y).
\]

\[
\begin{aligned}
\operatorname{diam}(A) &= \sup\{ d_X(x,y) : x,y \in A \} \\
\operatorname{diam}(f(A)) &= \sup\{ d_Y(f(x), f(y)) : x,y \in A \}.
\end{aligned}
\]

Since $d_Y(f(x),f(y)) = d_X(x,y)$ for all $x,y \in A$, we obtain
\[
\sup\{ d_Y(f(x), f(y)) : x,y \in A \}
=
\sup\{ d_X(x,y) : x,y \in A \}.
\]

Therefore
\[
\operatorname{diam}(f(A))=\operatorname{diam}(A).
\]
\end{proof}

\begin{remark*}[Dependencies]
\begin{itemize}
  \item Definition of isometry: $d_Y(f(x),f(y))=d_X(x,y)$ for all $x,y\in X$.
  \item Definition of image of a set: $f(A)=\{f(a):a\in A\}$.
  \item Definition of diameter: $\operatorname{diam}(A)=\sup\{d_X(x,y):x,y\in A\}$.
  \item Basic facts about $\sup$ and set inclusion (monotonicity of $\sup$).
\end{itemize}
\end{remark*}
